\chapter{Glossario dei termini Pāli}

\enlargethispage{2\baselineskip}

\begin{description}

\item[Anattā] Letteralmente, 'non-sé', cioè impersonale, senza essenza
  individuale; né una persona né appartenente a una persona. Una delle tre
  caratteristiche dei fenomeni condizionati.

\item[Anicca] Transitorio, impermanente, instabile, che ha la natura di
  sorgere e svanire. Una delle tre caratteristiche dei fenomeni
  condizionati.

\item[Añjali] Un gesto di rispetto. I palmi di entrambe le mani si uniscono
  direttamente davanti al petto, con le dita allineate e rivolte
  verso l'alto.

\item[Arahaṁ/Arahant] Letteralmente, 'degno' --- un termine applicato a tutti
  gli esseri illuminati. Come epiteto del solo Buddha, si usa 'Signore'.

\item[Ariyapuggalā] 'Esseri Nobili' o 'Nobili Discepoli' --- ne esistono
  otto tipi: coloro che stanno lavorando o che hanno raggiunto i quattro
  diversi stadi di realizzazione.

\item[Bhagavā] Generoso, fortunato --- quando usato come epiteto del
  Buddha, 'il Fortunato', 'il Beato'.

\item[Bhikkhu] Un monaco buddista che vive come mendicante, rispettando
  227 precetti di addestramento che definiscono una vita di rinuncia e
  semplicità.

\item[Brahmā] Essere celeste; un dio in uno dei regni spirituali
  superiori.

\item[Buddha] Colui che comprende, il Risvegliato, che conosce le cose
  così come sono; un potenziale in ogni essere umano. Il Buddha storico,
  Siddhattha Gotama, visse e insegnò in India nel V secolo a.C.

\item[Deva] Un essere celeste. Meno raffinato di un brahmā; poiché un deva si
  trova ancora in un regno sensuale, sebbene molto raffinato.

\item[Dhamma] (Sanscrito: Dharma) L'insegnamento del Buddha contenuto
  nelle scritture; non di carattere dogmatico, ma più simile a una zattera o
  un veicolo per condurre il discepolo alla liberazione. Anche, la Verità verso
  la quale quell'insegnamento punta; ciò che è al di là delle parole, dei concetti o
  della comprensione intellettuale. Quando scritto come '\emph{dhamma}', cioè
  con la 'd' minuscola, si riferisce a un 'elemento' o 'cosa'.

\item[Dukkha] Letteralmente, 'difficile da sopportare' --- disagio, irrequietezza
  della mente, angoscia, conflitto, insoddisfazione, scontento, stress,
  sofferenza. Una delle tre caratteristiche dei fenomeni condizionati.

\item[Fattori del Risveglio (bojjhaṅga)] 1.~consapevolezza, 2.~investigazione della verità, 3.~sforzo, 4.~rapimento, 5.~tranquillità, 6.~concentrazione, 7.~equanimità.

\item[Fondamenti della Consapevolezza (satipaṭṭhāna)] Consapevolezza di 1.~\emph{kāya} (corpo), 2.~\emph{vedanā} (sensazioni), 3.~\emph{citta} (mente),
4.~\emph{dhamma} (oggetti mentali).

\item[Tipi di Nascita (yoni)] Le quattro modalità di generazione con cui
  gli esseri nascono: nati da utero, nati da uovo, nati dall'umidità e nati spontaneamente.

\item[Vita Santa (brahmacariya)] Letteralmente: la condotta di Brahma; di solito
  riferito alla vita monastica. L'uso di questo termine enfatizza il voto di
  celibato.

\item[Jhāna] Assorbimento mentale. Uno stato di forte concentrazione focalizzato
  su una singola sensazione fisica o nozione mentale.

\item[Kamma] (Sanscrito: karma) Azione, atto; azioni create da
  impulso abituale, intenzione, volizione.

\item[Khandhā] I cinque aggregati, fisici o mentali ---
  cioè: \emph{rūpa, vedanā, saññā, saṅkhārā, viññāṇa.} L'attaccamento a
  uno qualsiasi di questi come 'Questo è mio', 'Io sono questo' o 'Questo è il mio sé' è
  \emph{upādāna} --- attaccamento o afferrare.

\item[Māra] Personificazione delle forze del male. Durante la lotta del Buddha
  per l'illuminazione, Māra manifestò forme spaventose e allettanti per
  tentare di distoglierlo dal suo obiettivo.

\item[Nibbāna] (Sanscrito: Nirvāṇa) Letteralmente, 'freschezza' --- lo stato di
  liberazione da ogni sofferenza e contaminazione, l'obiettivo del
  sentiero buddista.

\item[Paccekabuddha] Buddha solitario --- qualcuno illuminato con i propri
  sforzi senza fare affidamento su un insegnante ma che, a differenza del Buddha, non ha
  un seguito di discepoli.

\item[Paritta] Versi cantati in particolare per benedizione e protezione.

\item[Parinibbāna] Il passaggio finale del Buddha, cioè l'ingresso finale nel
  Nibbāna.

\item[Saggio Pacifico (muni)] Un epiteto del Buddha.

\item[Piani di Nascita (bhūmi)] I tre piani in cui avviene la rinascita:
  \emph{kāmāvacara-bhūmi}: il piano sensuale;
  \emph{rūpāvacara-bhūmi}: il piano della forma; \emph{arūpāvacara-bhūmi}: il piano
  senza forma.

\item[Puñña] Merito, l'accumulo di buona fortuna, benedizioni o
  benessere risultante dalla pratica del Dhamma.

\item[Rūpa] Forma o materia. Gli elementi fisici che compongono il corpo,
  cioè terra, acqua, fuoco e aria (solidità, coesione, temperatura e
  vibrazione).

\item[Saṅgha] La comunità di coloro che praticano la Via del Buddha.

  Più specificamente, coloro che si sono formalmente impegnati nello stile di vita
  di monaci e monache mendicanti. Le 'quattro coppie, gli otto tipi di
  esseri nobili' sono coloro che sono sul sentiero o che hanno realizzato la
  fruizione dei quattro stadi dell'illuminazione: entrata nella corrente, una sola volta ritornante,
  non-ritornante e arahant.

\item[Saṅkhārā] Formazioni, costruzioni, tutte le cose condizionate, o impulsi
  volitivi, cioè tutti gli stati mentali eccetto la sensazione e la percezione
  che colorano i propri pensieri e li rendono buoni, cattivi o neutri.

\item[Saññā] Percezione, la funzione mentale del riconoscimento.

\item[Tathāgata] 'Così andato' o 'Così venuto' --- colui che è andato oltre
  la sofferenza e la mortalità; colui che sperimenta le cose come sono,
  senza illusione. L'epiteto che il Buddha applicava a se stesso.

\item[Triplice beatitudine] Beatitudine mondana, beatitudine celeste e beatitudine
  del Nibbāna.

\item[Triplice Gemma] Buddha, Dhamma e Saṅgha.

\item[Vedanā] Sensazione --- sensazioni fisiche e mentali che possono essere
  piacevoli, spiacevoli o neutre.

\item[Viññāṇa] Coscienza sensoriale --- il processo attraverso il quale si ha
  il vedere, l'udire, l'odorare, il gustare, il toccare e il pensare.

\end{description}
